\documentclass{article}
\usepackage{amsmath,amssymb}
\usepackage[margin=1in]{geometry}
\allowdisplaybreaks
\usepackage[UTF8]{ctex}
\begin{document}

\section*{集成式的交通感知动态调度系统，以及闲置车辆的重新定位功能}

\subsection*{派遣目标与重新安置目标，以及与之相关的主体-区块信息}
继承的派发规则是：
\begin{equation}
\{v^*,x_t^{v^*}\}=\operatorname*{argmin}_{v,x}
\bigl[c(v,x_t^v)-c(v,x)\bigr].\tag{1}
\end{equation}
搜索范围包括几乎所有附近的车辆，以及在这些车辆中添加新请求后所更新的路线。其路线成本如下：
\begin{equation}
c(v,x)=\gamma T(v,x)+(1-\gamma)
\left[\beta T(v,x)^2+\sum_{n\in P_v}Y_n(v,x)\right].\tag{2}
\end{equation}
在迁移时代$h$时，迁移的目标为：
\begin{equation}
\Phi=\min\;\sum_{i\in\bar N}\sum_{j\in\bar N}\lambda_i t_{ij}X_{ij}
+\theta\sum_{i\in\bar N}\sum_{j\in\bar N}r_{ij}W_{ij}.\tag{3}
\end{equation}
受以下条件限制：
\begin{align}
\sum_{j\in\bar N}X_{ij}&=1,
&&\forall i\in\bar N,\tag{4}\\
Y_{jm}&\leq Y_{j,m-1},
&&\forall j\in\bar N,\quad m=2,3,\ldots,C_j,\tag{5}\\
\sum_{i\in\bar N}\lambda_iX_{ij}
&\leq\mu_j\left[Y_{j1}\rho_{\eta j1}
+\sum_{m=2}^{C_j}Y_{jm}(\rho_{\eta jm}-\rho_{\eta j,m-1})\right],
&&\forall j\in\bar N,\tag{6}\\
X_{ij}&\leq Y_{j1},&&\forall i,j\in\bar N,\tag{7}\\
\sum_{j\in\bar N}\sum_{m=1}^{C_j}Y_{jm}&=B,\tag{8}\\
\sum_{j\in\bar N}W_{ij}&=S_i,&&\forall i\in\bar N,\tag{9}\\
\sum_{i\in\bar N}W_{ij}&=D_j,&&\forall j\in\bar N,\tag{10}\\
S_j&\leq y_j,&&\forall j\in\bar N,\tag{11}\\
-D_j+S_j-y_j+\sum_{m=1}^{C_j}Y_{jm}&\leq0,
&&\forall j\in\bar N,\tag{12}\\
X_{ij}&\in\{0,1\},\tag{13}\\
0\leq Y_{jm}&\leq1,\tag{14}\\
D_j,S_j&\geq0,\qquad W_{ij}\in Z^+.\tag{15}
\end{align}
这两个目标都是事件驱动型运营策略的组成部分，而不是单一总体目标的条款。下面的“ transit-aware dispatch”结构属于P1；而“relocation”则属于P2。方程(11)到(12)是作者提出的替代不等式，它们考虑了节点处的流入和流出情况。被替代的流量守恒方程则没有被强制应用。

\subsection*{决策角色、固定输入以及路由条件}
在 (1)–(2) 中，$v$代表车辆，$x$代表其行程；$x_t^v$是在插入请求后新增的行程。$T(v,x)$表示行程所需的时间长度。$P_v$表示分配给车辆$v$的乘客数量，而$Y_n(v,x)$则是乘客$n$的旅行时间，即等待时间加上在车内的旅行时间。这个乘客时间函数与重新定位变量$Y_{jm}$是不同的。系数$\gamma$用于将客户成本与系统成本进行转换，而$\beta$则是预测程度参数。调度队列延迟的近似值被参数化为$\beta T(v,x)^2$；第4页上打印的M/M/1表达式则……
\[
\frac{\mu\lambda}{2(\mu-\lambda)}T(v,x)^2.
\]
这并不表示$\beta$等于那个系数。打印出来的边界情况实际上是$\beta=0$，对于这种情况，调度行为完全表现为近视行为；它并不是所选择的非近视版本。对于$\beta$或$\gamma$来说，这里并没有固定的数值。

路由网络就是完整的图。
\[
G(N,E),\qquad N=N_T\cup N_P\cup N_Z,
\qquad V=\{v_1,v_2,\ldots,v_{|V|}\}.
\]
$N_T$是中转站的集合，$N_P$是请求中的取送点，而$N_Z$则是车辆可以停放的区域中心点。$E$是连接点集合。在路线规划方面，$t_{ij}$是从节点$i$到节点$j$的最短路径行驶时间。$V$是一个具有固定容量的车队，而$P$则是所有乘客的集合。一个请求会指定其取送地点、取送时间等参数。每个节点$i\in N$的到达情况都假设遵循泊松过程，其到达率為$\lambda_i$。

Tour的构建必须满足“先接载再送离”的优先顺序，同时还要考虑车辆的承载能力。本文通过旅行推销员问题来求得$T(v,x)$和$x$的值。文中提供了一种插入式启发式算法，而不是单独的代数可行集公式化方法。没有对请求时间窗口进行限制。所有请求都能得到处理，用户在等待期间也不会取消请求。这些假设在保留的情况下被保留下来，没有增加拒绝阈值或等待时间限制。

在通勤途中，最多只能进行两次跨模式的换乘。不允许在两种不同的共享出行车辆之间进行直接换乘。每个通勤前的或通勤后的路段，都由一辆共享出行车辆或步行来完成。通勤时间的估算是基于固定的公共交通时间表。共享出行车辆的容量没有限制，前提是假设共享出行用户只是普通通勤乘客中的一小部分。

服务区域$\mathcal A$被划分为一系列搬迁区域$\bar N$。搬迁过程采用区域间传输方式$t_{ij}$，这与节点间路由的传输时间有所不同。该过程的固定输入在某一时间点为：
\begin{description}
\item[$\lambda_i$] 表2中，在前一个迁移时期$h-1$内，区域$i$的到达率；P2使用预测的时期$h$的到达率，包括预期的出站车站到达情况，具体描述如下。
\item$\mu_j$区域$j$在之前的搬迁时期$h-1$的服务率，如表2所示；综合策略会实时更新服务率估算值。
\item$r_{ij}$从区域$i$迁移到区域$j$的迁移成本。
\item$B$在纪元开始时的闲置车辆总数$h$。
\item$C_j$是区域$j$中最大可能的闲置车辆数量。
\item$y_j$在纪元开始时的区域$j$中，闲置车辆的数量为$h$。
\item$\theta$在重新定位目标中的转换标量。
\item[$\rho_{\eta jm}$] 在区域$j$中，针对$\eta$和$m$服务器的可靠性，需要确定外部队列强度系数。表2列出了符号$\rho_{\eta m}$；(6)中包含了区域索引。
\end{description}
在(3)--(15)中，时代索引被省略了。搬迁决策包括：$W_{ij}$，即闲置车辆从区域$i$转移到区域$j$；$X_{ij}$，当客户在区域$i$到达时，从区域$j$提供服务；$Y_{jm}$，表示在区域$j$有第$m$辆闲置车辆来为顾客服务，同时持续进行放松操作(14)；还有虚拟的供需变量$S_i,D_j$。它们的区域索引范围在$\bar N$之间，而车辆位置索引则从$1$到$C_j$。方程(13)--(15)以原始顺序打印，没有重复的量词；$Z^+$则保持原样。

区域$i$内的所有客户都被分配到同一个服务区域$j$。 方程（5）用于排列车辆的位置。 方程（8）将他们整个系统的总和固定为整数$B$。 作者通过分段线性队列约束和这一整数总量，来证明将$Y_{jm}$从二进制形式转换为连续形式在(14)中的合理性；没有添加额外的区域整体性条件。 选择保留（6），而不是那种过于短视的替代方案，后者会放宽这一要求。 附录中描述了这种重新定位的近似模型，其形式为 M/M/$s$，并且服务时间呈指数分布。

\subsection*{外源队列强度构建}
对于用户定义的可靠性$\eta$、服务器数量$m$，以及其他等待的客户数量$b$每一种组合，论文建议通过求解方程（16）的根来间接确定强度系数，然后将该值代入方程（6）中。
\begin{equation}
\sum_{k=0}^{m-1}
\left(\frac{(m-k)m!\,m^b}{k!}\right)
\left(\frac{1}{\rho^{m+b+1-k}}\right)
\geq\frac{1}{1-\eta}.\tag{16}
\end{equation}
这里的$k$是求和指标，而不是P1所使用的最近站点的数量。式(6)的含义是，不超过$b$的其他顾客等待的时间，其概率应大于服务可靠性$\eta$。

\paragraph{印刷资格。}
尽管附带的文字描述为“寻找根”，但(16)被打印为不等式形式。所显示的不等式被保留了下来。在操作过程中，正 equality边界被用作系数；论文中描述了一个根，而作者提供的40个公共系数表中的所有值，在四舍五入后都符合这一边界。这是一个参数构造关系，而不是对重新定位决策变量的约束。关于Marianov和Serra在2002年提出的进一步推导内容并未提供，因此仍然未被验证。实验中的选择$\eta=0.95$和$b=0$并不是这种公式化的通用输入限制条件。

\subsection*{P1：具有交通感知能力的调度和路线更新}
当收到来自$o$向$d$的新请求$n$时，需要更新从请求时间$n-1$开始以来每辆车辆的当前位置和服务状态。对于附近的多辆车辆，需要考虑以下几种情况：仅拼车$r$、拼车+转运+拼车$rtr$，以及拼车+转运+步行或其反向流程$rtw$、$wtr$。论文中用相应的符号来表示这些成本。
\[
C_r^n,\qquad C_{rtr}^n,\qquad C_{rtw}^n.
\]
对于交通出行方案，需要考察每对$(s_1,s_2)$中的$k$最近的入口车站与$o$和$k$最近的输出车站与$d$之间的位置关系。 多模式旅行时间是指从出发地到进入特定模式的旅行时间、在车站的换乘等待时间、在途中的旅行时间、在出口车站的等待时间，以及从出口车站到目的地特定模式的旅行时间的总和。 第二个拼车段是根据当前系统状态进行近似计算的。P1建议选择这些站点对之间最短的多模式出行时间，并计算出最快的出行方案以及相应的车辆分配。

在 P1 中使用的运输成本是
\[
\text{transit cost}
=\tfrac12\,\text{headway}+\text{travel time in the transit system}.
\]
P1，步骤2(a)，打印出以下成本最小化指令：
\begin{itemize}
\item 对于$r$，使用公式(2)来求出$(o,d)$的值，这样……
\[
[c(v,x)-c(v,x_t^v)]
\]
被最小化。
\item 对于$rtr$，使用公式(2)来分别求解每个$(o,s_1)$和$(s_2,d)$，使得它们的总和等于…
\[
[c(v,x)-c(v,x_t^v)]
\]
对于每一次出行，加上交通费用后，总成本都能降到最低。
\item 对于$rtw$（$wtr$），使用公式（2）来求解每一个$(o,s_1)$（或$(s_2,d)$），从而…
\[
[c(v,x)-c(v,x_t^v)]
\]
对于拼车出行来说，加上交通费和步行离开的费用，总成本被降到最低。
\end{itemize}
在选择了服务选项之后，当请求涉及使用公共交通方式时，请更新取送点的位置信息；同时，也要更新分配给该车辆的路线信息。

\paragraph{印刷分发不一致。}
第9页的P1最小化了$c(v,x)-c(v,x_t^v)$，即(1)中插入成本表达式的负值，而$x_t^v$作为更新后的行程定义保持不变。这两个打印符号分别进行评估。主要的操作选择是(1)中的正更新-负旧成本，这与原始文本中最小化额外插入成本的观点一致。同时，也报告了P1的负成本选项。论文还描述了“最快”的选择，同时规定了这些成本最小化方案。没有引入替代选择标准。第8页关于获取$T(v,x)$和$x$“用于方程(3)”的引用指向一个不包含这些符号的方程；它们的打印路线成本定义为(2)。

\subsection*{P2和迁移接口}
当车辆完成服务后，它就会处于闲置状态。在迁移时期$h$开始时，P2会识别出处于闲置状态的车辆及其当前所在的区域，并预测$\lambda$在时期$h$内的情况，包括预计需要共享出行服务的车辆到达中转站的时间。它会使用方程(3)--(15)来解决问题，其中外生系数来自方程(16)。如果一辆闲置车辆被分配到不同的区域，那么该区域的新位置会被添加到该车辆的行程中。

一个$rtr$的选择并不会立即将车辆安排进行过境运输。P1包含了预期的拼车费用，而这一选择则会增加在出发站的需求，尤其是在预期的到达时间时。当乘客到达出发站后，会有一列仅包含拼车的车辆将乘客送往最终目的地。因此，运输安排是根据乘客的到达时间来进行的，而重新安排则根据预定义的时刻进行；这两种时间安排是相互连接的。

在第11至12页的实现部分中，会更新车辆和乘客的信息，使其与当前的时间同步。当乘客在进站站被放下时，系统会计算等待下一列火车的時間。对于$rtr$乘客来说，系统会生成一个新的请求，将出站站作为接送点，最终目的地作为下车点，并将期望的接送时间设置为预计到达出站站的时间。该请求会被添加到请求列表中，该列表按照接送时间进行排序。而$rtw$乘客则乘坐下一列火车，然后步行前往目的地。此时不会立即为后续车辆进行预订。

所选的途中切换版本允许在搬迁过程中的闲置车辆接载新客户，这与报告中的集成策略实验以及附录4中的内容一致。当一辆闲置车辆到达其指定的搬迁地点时，它的状态就会变得可用，然后它就会从过渡期间的闲置车辆列表中移除。对于搬迁计算来说，该实现方式会更新闲置车辆列表，同时排除那些仍处于过渡阶段的车辆。

在搬迁间隔长度为$\Delta$、预热时间为$T_{\mathrm{warm}}$、当前模拟时间为$t$，以及纪元计数器为$h$的情况下，该实现将进行初始化。
\[
t=0,\qquad h=0,\qquad
\mu_i=\mu_0\quad\text{for each zone }i=1,\ldots,\bar N.
\]
在更新状态并根据请求进行调度之后，其打印出的重新定位触发器和时代更新信息会随之更新。
\[
t>\max(h\Delta,T_{\mathrm{warm}}),\qquad h:=h+1.
\]
当这一严格触发条件成立时，到达率$\lambda_i^h$、服务率$\mu_i^h$以及区域中心位置$(x_i^{*h},y_i^{*h})$将通过三步移动平均法进行更新。P2会确定$W_{ij}$的值。请求会被处理，直到所有新请求都得到服务为止；之后，时钟会继续推进，处理剩余的行程，直到所有客户都得到服务。这里没有固定的实例特定间隔、预热时间、初始仓库位置或车队规模。

\subsection*{在线价格以及动态更新的搬迁中心点}
对于$i\in N_Z$，所实现的平均服务速率是在一个时间间隔内定义的。
\[
\mu_i=
\frac{\text{total customers served with pickup in zone }i}
{\text{total in-vehicle time of those served customers}}.
\]
该政策采用三步移动平均法来学习服务费率，以避免在搬迁间隔期间出现波动。在实施过程中，到达率、服务费率以及质心更新值均基于前三个搬迁时期的均值进行计算。服务费率取决于调度和路由策略、客户到达率以及车辆位置。这些费率是P2模型的估算输入，并非额外的优化变量。

需求加权重心计算是
\begin{equation}
x_i^*=\frac{\sum_{s=1}^{m}\lambda_s x_s}{\sum_{s=1}^{m}\lambda_s},
\qquad
y_i^*=\frac{\sum_{s=1}^{m}\lambda_s y_s}{\sum_{s=1}^{m}\lambda_s},
\qquad\forall s\in i,\quad\text{for }i=1,\ldots,\bar N.
\tag{17}
\end{equation}
这里的$x_s,y_s$表示在区域$i$内，从$s$点出发的取货点的坐标；而$\lambda_s$则表示到达$s$点的速率。用于重新定位的区域的质心，与用于取货和送货的顾客起点和终点位置是分开的。这些质心会在每个时期进行更新，以反映顾客需求的时空分布情况。

\paragraph{印刷范围与特殊情况限制。}
方程（17）保留了打印出的范围$\forall s\in i$和范围$i=1,\ldots,\bar N$。不过，表2将$\bar N$定义为一个集合，而非基数。其上限值$m$在取点计算中有所使用，但在相关段落中并未给出单独的本地定义；它也没有与(5)--(6)和(16)中的服务器索引相对应。现有的公式并未为(17)或服务速率比提供零分母规则，也没有为不完整的三个时期的平均历史数据提供启动规则。因此，没有添加任何内容。

\paragraph{来源范围。}
采用 Ma、Rasulkhani、Chow 和 Klein 的论文 \emph{融合公交换乘的动态拼车派单与空车重定位策略}（DOI: 10.1016/j.tre.2019.07.002）PDF 实际第 4 至 12 页所列的模型。排队假设和选定的途中派单策略依据所提供的附录。式 (16) 至 (17) 及 P1/P2 均来自所提供的正文。未提供正文引用的既有派单、排队推导和行程构建文献，此处未引入其中的额外模型。实验标定、用于比较的空车重定位策略及其目标在所选模型的范围之外。


\paragraph{有限的执行证据。}
作者的公开表格\texttt{rho\_0\_95\_m\_40\_b0.txt}位于仓库中
\texttt{BUILTNYU/data\_双模态\_拼车}，提交
\texttt{303f8806add13cdc7390a0efee268be0eb7c0d2a}，确认了所有四十个服务器计数的正根约定，在可靠性\(0.95\)下
。提供的输入保持其自身的\(\eta=1/2\)，并使用合理的隔离间隔来计算相应的根。既评估了\(W_{ij}\ge0\)整数根，也评估了\(W_{ij}\ge1\)整数根。请求到达量来自三个外部接纳周期上的两个标记的泊松过程。严格的事件触发可能在普通接纳结束后再次触发一个引发的请求。物理上完成的有限轨迹报告了阶段最优的搬迁和启发式源调度，而不是不存在的集中最优解。对于步行-交通-拼车，运营商会立即预订唯一的车辆，乘客不能在实际离开的火车到达之前上车。车辆可能会在此发布之后进一步插入。对于拼车-交通-拼车，只有第一辆车在初始调度时被预订；第二辆车在火车离开时调度。步行到达是实际的最终事件。
\end{document}

